% !PRINTED 39
\weberpage{39}
Zugleich ist, wie die Formel (17) zeigt, wenn $\Theta_1$ ein beliebiges Element in $S$ ist, $\Theta_1^{a_\nu}$ in $S_{\nu-1}$ enthalten, und die Exponenten $\lambda_1, \lambda_2 \ldots \lambda_{\nu-1}$ sind durch $a_\nu$ theilbar. Es ist daher auch die Forderung 2) befriedigt.

Damit ist also unser Satz I. bewiesen. Zugleich sehen wir aus dieser Ableitung, \gesperrt{dass man für eine beliebige Abel'sche Gruppe $S$ die Elemente der Basis immer so annehmen kann, dass ihre Grade Primzahlpotenzen sind.}

Aus der Darstellbarkeit durch eine Basis folgt auch, dass jede Abel'sche Gruppe metacyklisch ist; denn sind die Elemente einer solchen Gruppe $S$ in der Form dargestellt:
\[
A = A_1^{\alpha_1} A_2^{\alpha_2} \ldots A_\nu^{\alpha_\nu},
\]
und ist $p$ eine im Grade $a_1$ von $A_1$ aufgehende Primzahl, so bilden die Elemente
\[
A' = A_1^{\alpha_1 p} A_2^{\alpha_2} \ldots A_\nu^{\alpha_\nu},
\]
wenn $\alpha_1$ ein volles Restsystem nach dem Modul $a_1 : p$ durchläuft, einen Theiler $S'$ von $S$ vom Index $p$, der durch die Basis $A_1^p, A_2 \ldots A_\nu$ darstellbar ist. Von $S'$ kann man wieder auf die gleiche Weise einen Theiler von Primzahlindex finden u. s. f. Also ist $S$ metacyklisch (\S.\,8).

\begin{center}
\S.\,10.

\gesperrt{Die Invarianten der Abel'schen Gruppen.}
\end{center}

Der Beweis, den wir im vorigen Paragraphen für die Möglichkeit der Darstellung einer Abel'schen Gruppe durch eine Basis mitgetheilt haben, enthält zugleich einen Weg, eine solche Basis zu finden, und zwar eine, bei der die Grade der Basiselemente Primzahlpotenzen sind. Gleichwohl kann es vorkommen, dass eine und dieselbe Gruppe auf verschiedene Arten durch Basen dargestellt werden kann.

Betrachten wir z. B. zwei Elemente $A, B$, deren Grade $a, b$ relativ prim sind, so wird durch diese als Elemente einer zweigliedrigen Basis eine Gruppe
\weq{1}{A^\alpha B^\beta \qquad \begin{aligned}
\alpha &= 0, 1, 2 \ldots a-1 \\
\beta &= 0, 1, 2 \ldots b-1
\end{aligned}}
dargestellt. Dieselbe Gruppe kann aber auch durch die eingliedrige Basis $AB$ dargestellt werden in der Form
\weq{2}{(AB)^s \qquad s = 0, 1, 2 \ldots ab-1.}
